\scanpage{6}
de $\underline O_X$ sur les $Q^i=\mathcal P^{i+1}$ \unsure{commutent} à cette structure. Il s'ensuit que les
$Q^i[M]=\underset{\inserted{\text{[str. ext. dr.]}}}{Q^i}\otimes_{\underline O_X}M$
sont les composantes d'un objet \deleted{\unsure{semi-}}cosimplicial, noté $Q^*[M]$, qui est en fait un \uline{Module} cosimplicial sur l'anneau cosimplicial $\mathcal P^*$.
Ce $Q^*[M]$ dépend également de façon fonctorielle de $M$, quand $M$ parcourt $\operatorname{Mod}(X/S)$.

3. Notons un caractère hautement spécial des Modules cosimpliciaux de la forme $Q^*[M]$ sur $\mathcal P^*$ : \deleted{ils} \deleted{\unsure{sont des}} c'est que les $Q^i[M]$ se déduisent les uns des autres par changement de base via les opérations simpliciales $\mathcal P^i\longrightarrow\mathcal P^j$ [i.e. ce sont des \deleted{\unsure{petits}} objets d'un $\varprojlim$, non seulement
\margin{5}
